\section*{Bab \ref{ch:g11:gene-expression} --- Dari Gen ke Protein}

\begin{solution}{exo:g11:gene-expression:1}
\omterm{def:g11:gene-expression:transcription}{RNA} berutai tunggal, gulanya ribosa, dan ia memakai urasil sebagai
ganti timin (ia juga berumur pendek dan meninggalkan \omterm{def:g10:cells-common-unit:organelle}{inti selnya}).
\end{solution}

\begin{solution}{exo:g11:gene-expression:2}
\texttt{5'-AUG GCU UUU UAA-3'}.
\end{solution}

\begin{solution}{exo:g11:gene-expression:3}
Met--Gly--Lys--Cys, lalu henti.
\end{solution}

\begin{solution}{exo:g11:gene-expression:4}
\omterm{def:g10:cells-common-unit:organelle}{Ribosomnya} membaca mRNA \omterm{def:g11:gene-expression:code}{kodon} demi \omterm{def:g11:gene-expression:code}{kodon} dan menyambungkan \omterm{def:g10:chemistry-of-life:families}{asam
aminonya} menjadi rantai; tiap \omterm{def:g11:gene-expression:transcription}{RNA} transfer mencocokkan satu \omterm{def:g11:gene-expression:code}{kodon}
dengan antikodonnya dan membawa \omterm{def:g10:chemistry-of-life:families}{asam amino} yang sepadan.
\end{solution}

\begin{solution}{exo:g11:gene-expression:5}
Dengan 4 huruf, kata berhuruf dua memberi $4^2 = 16$ gabungan, lebih
sedikit daripada 20 \omterm{def:g10:chemistry-of-life:families}{asam amino}; kata berhuruf tiga memberi 64, dan itu cukup.
\end{solution}

\begin{solution}{exo:g11:gene-expression:6}
$412 \times 3 + 3 = 1239$ \omterm{def:g10:universal-dna:nucleotide}{nukleotida}; sebuah \omterm{def:g10:universal-dna:gene}{gen} sepanjang sedikitnya
1239 pasangan \omterm{def:g10:universal-dna:nucleotide}{basa} (lebih panjang pada eukariot, dengan \omterm{prop:g11:gene-expression:splicing}{intronnya}).
\end{solution}

\begin{solution}{exo:g11:gene-expression:7}
Aslinya: Met--Lys--Gly--Trp. Mutannya: UGA adalah henti, sehingga
Met--Lys--Gly: \omterm{def:g11:mutations:mutation}{mutasi} tanpa makna, yang memenggal \omterm{def:g11:gene-expression:protein}{proteinnya}.
\end{solution}

\begin{solution}{exo:g11:gene-expression:8}
\texttt{AAG} tetap Lys: diam, Met--Lys--Gly--Trp tidak berubah. Yang
kedua adalah insersi G yang menggeser rangkanya: \texttt{AUG AAA GGG
CUG GUA A}\dots{} Met--Lys--Gly--Leu--Val\dots, sebuah pergeseran
rangka tanpa henti dalam penggalannya.
\end{solution}

\begin{solution}{exo:g11:gene-expression:9}
Pada sebagian besar blok tabelnya, keempat \omterm{def:g11:gene-expression:code}{kodon} yang berbagi dua huruf
pertama menetapkan \omterm{def:g10:chemistry-of-life:families}{asam amino} yang sama (Pro, Thr, Ala, Gly, Val, Ser,
Leu\dots): mengubah huruf ketiganya lalu tidak mengubah apa pun.
Degenerasi kodenya terutama terletak pada kedudukan ketiganya.
\end{solution}

\begin{solution}{exo:g11:gene-expression:10}
Jika dibaca bertiga, \texttt{UCU CUC UCU CUC}\dots{} berselang-seling
dua \omterm{def:g11:gene-expression:code}{kodon}, yaitu UCU dan CUC, sehingga \omterm{def:g11:gene-expression:protein}{proteinnya} berselang-seling dua
\omterm{def:g10:chemistry-of-life:families}{asam amino} --- seperti yang teramati. Dengan kode berhuruf dua,
pengulangan yang sama akan memberi satu \omterm{def:g11:gene-expression:code}{kodon} dan satu \omterm{def:g10:chemistry-of-life:families}{asam amino}.
Karena poli-U memberi Phe dan pada tabelnya UCU adalah Ser dan CUC
adalah Leu, percobaannya menetapkan UCU sebagai Ser dan CUC sebagai Leu
(atau sebaliknya, yang diputuskan percobaan lain).
\end{solution}

\begin{solution}{exo:g11:gene-expression:11}
Hanya 1503 pasangan \omterm{def:g10:universal-dna:nucleotide}{basa} yang diperlukan bagi urutan penyandinya;
sisanya berupa \omterm{prop:g11:gene-expression:splicing}{intron}, yang ditranskripsikan lalu dipotong keluar dari
pra-mRNA-nya, ditambah urutan pengatur pada ujung \omterm{def:g10:universal-dna:gene}{gennya}.
\end{solution}

\begin{solution}{exo:g11:gene-expression:12}
Menghalangi translasi membunuh bakterinya sambil membiarkan \omterm{def:g10:cells-common-unit:cell}{sel}
pasiennya tetap bekerja. Kekhususan obatnya bermakna bahwa \omterm{def:g10:cells-common-unit:organelle}{ribosom}
bakteri dan \omterm{def:g10:cells-common-unit:organelle}{ribosom} manusia, walaupun mengerjakan tugas yang sama
dengan kode yang sama, cukup berbeda strukturnya sehingga sebuah
molekul dapat mengikat yang satu dan tidak yang lain --- yaitu beda
yang menumpuk sejak keduanya berpisah.
\end{solution}

\begin{solution}{exo:g11:gene-expression:13}
\omterm{def:g11:gene-expression:transcription}{RNA} polimerase mencitnya menstranskripsikan \omterm{def:g10:universal-dna:gene}{gen} ubur-ubur itu, dan
\omterm{def:g10:cells-common-unit:organelle}{ribosom} serta tRNA-nya menerjemahkan mRNA-nya dengan tabel \omterm{def:g11:gene-expression:code}{kodon} yang
sama, sehingga urutan \omterm{def:g10:chemistry-of-life:families}{asam amino} yang sama, dan dengan demikian \omterm{def:g11:gene-expression:protein}{protein}
berpendar yang sama, pun terbuat. Dengan kode yang berbeda, \omterm{def:g11:gene-expression:code}{kodonnya}
akan dibaca sebagai \omterm{def:g10:chemistry-of-life:families}{asam amino} lain dan \omterm{def:g11:gene-expression:protein}{proteinnya} akan berupa rantai
tanpa makna yang tidak berpendar.
\end{solution}

\begin{solution}{exo:g11:gene-expression:14}
UAG tidak lagi akan menghentikan translasinya: setiap \omterm{def:g11:gene-expression:protein}{protein} yang
\omterm{def:g10:universal-dna:gene}{gennya} berakhir pada UAG akan diperpanjang melampaui ujung normalnya,
dan sering kehilangan fungsinya --- sebuah kerugian yang meluas. Tetapi
\omterm{def:g10:universal-dna:gene}{gen} yang terpenggal oleh UAG yang terlalu dini kini akan dibaca terus,
sehingga \omterm{def:g11:gene-expression:protein}{protein} yang nyaris utuh panjangnya pulih: \omterm{def:g11:mutations:mutation}{mutasi} kedua
menekan \omterm{def:g11:mutations:mutation}{mutasi} pertama.
\end{solution}

\begin{solution}{exo:g11:gene-expression:15}
Pergeseran rangka mengacaukan segalanya di hilirnya dan lazimnya
bertemu sebuah henti dalam beberapa \omterm{def:g11:gene-expression:code}{kodon}: di dekat awalnya, pada
dasarnya tidak ada \omterm{def:g11:gene-expression:protein}{protein} yang benar yang terbuat; di dekat ujungnya,
sebagian besar \omterm{def:g11:gene-expression:protein}{proteinnya} utuh dan mungkin masih dapat melipat dan
bekerja. \omterm{def:g11:mutations:mutation}{Mutasi} salah makna mengubah satu \omterm{def:g10:chemistry-of-life:families}{asam amino}: tak berbahaya
jika kedudukan itu menahannya, mematikan jika kedudukan itu berada pada
tapak aktifnya atau merusak lipatannya --- pengaruhnya bergantung pada
\omterm{def:g10:chemistry-of-life:families}{asam amino} yang mana, di mana.
\end{solution}

\begin{solution}{pb:g11:gene-expression:1}
\textbf{1.} \texttt{3'-TAC CGA ACC TTT GGG
CTT ATG CCA GTA AAG ACT-5'}.

\textbf{2.} \texttt{AUG GCU UGG AAA CCC
GAA UAC GGU CAU UUC UGA}.

\textbf{3.} Met--Ala--Trp--Lys--Pro--Glu--Tyr--Gly--His--Phe, lalu henti.

\textbf{4.} 10 \omterm{def:g10:chemistry-of-life:families}{asam amino}; 33 \omterm{def:g10:universal-dna:nucleotide}{nukleotida} (11 \omterm{def:g11:gene-expression:code}{kodon} termasuk hentinya).

\textbf{5.} AUG; kedudukan AUG yang pertama itu menetapkan rangkanya,
dan tiap \omterm{def:g11:gene-expression:code}{kodon} berikutnya dibaca seirama darinya.

\textbf{6.} CCG tetap Pro: diam, \omterm{def:g11:gene-expression:protein}{proteinnya} tidak berubah.

\textbf{7.} CGU adalah Arg sebagai ganti His: salah makna,
Met--Ala--Trp--Lys--Pro--Glu--Tyr--Gly--Arg--Phe.

\textbf{8.} UGA adalah henti: Met--Ala, tanpa makna. Satu substitusi G
menjadi A pada \omterm{def:g11:gene-expression:code}{kodon} ketiganya (TGG menjadi TGA).

\textbf{9.} mRNA-nya \texttt{AUG GCU UUG GAA ACC
CGA AUA CGG UCA UUU CUG A}:
Met--Ala--Leu--Glu--Thr--Arg--Ile--Arg--Ser--Phe--Leu\dots, tanpa henti
dalam penggalannya: pergeseran rangka, dengan tiap \omterm{def:g10:chemistry-of-life:families}{asam amino} setelah
yang kedua berubah.

\textbf{10.} Mutan 1 (diam, tanpa perubahan) $<$ mutan 2 (satu \omterm{def:g10:chemistry-of-life:families}{asam
amino} berubah, fungsinya barangkali tetap) $<$ mutan 4 (pergeseran
rangka, semua \omterm{def:g10:chemistry-of-life:families}{asam aminonya} keliru kecuali dua) $\approx$ mutan 3
(tanpa makna, sebuah penggalan berisi dua \omterm{def:g10:chemistry-of-life:families}{asam amino}: \omterm{def:g11:gene-expression:protein}{proteinnya} lenyap).

\textbf{11.} Leu 6, Ser 6, Trp 1. Trp, yang berkodon tunggal, diubah
oleh substitusi apa pun; Leu atau Ser, yang berkodon enam, paling
sering tidak berubah.

\textbf{12.} CUU, CUC, CUA, CUG semuanya Leu: setiap satu dari 3
substitusi yang mungkin pada kedudukan ketiganya bersifat diam --- 100\%.

\textbf{13.} Dua huruf memberi 16 \omterm{def:g11:gene-expression:code}{kodon}, terlalu sedikit untuk 20 \omterm{def:g10:chemistry-of-life:families}{asam
amino} ditambah sebuah henti; tiga huruf memberi 64, yang sudah lebih
dari cukup, sehingga empat huruf (256) akan mubazir.

\textbf{14.} UUU $=$ Phe dan AAA $=$ Lys.

\textbf{15.} \omterm{def:g10:universal-dna:gene}{Gen} yang dipindahkan antarspesies dibaca serupa, sehingga
\omterm{prop:g10:universal-dna:universal}{transgenesis} berhasil; dan kode bersama, yang diwariskan tanpa berubah,
menjadi bukti bahwa seluruh organisme berasal dari nenek moyang bersama
yang sudah memakainya.

\textbf{16.} 10 \omterm{def:g10:chemistry-of-life:families}{asam amino}: \qty{1}{s}. 600 \omterm{def:g10:chemistry-of-life:families}{asam amino}: \qty{60}{s}.

\textbf{17.} $1800/80 \approx 22$ \omterm{def:g10:cells-common-unit:organelle}{ribosom} sekaligus.

\textbf{18.} $7200/8 = 900$ salinan.

\textbf{19.} Satu mRNA memberi $600/8 = 75$ salinan dalam 10 menit;
sekitar 27 mRNA diperlukan.

\textbf{20.} \omterm{def:g11:gene-expression:protein}{Proteinnya}
Met--Ala--Trp--Lys--Pro--Glu--Tyr--Gly--His--Phe; mutan 3, yaitu satu
substitusi G menjadi A yang menciptakan henti pada \omterm{def:g11:gene-expression:code}{kodon} ketiganya,
yang memusnahkannya; satu \omterm{def:g10:cells-common-unit:organelle}{ribosom} membangun kesepuluh \omterm{def:g10:chemistry-of-life:families}{asam aminonya}
dalam sekitar satu sekon.
\end{solution}
